\section{Appendix: Faithfulness Theorem -- Full Proof}\label{app:proof}

We restate \cref{thm:faithfulness} for self-containment.

\paragraph{Restated theorem.}
Let \(\varepsilon_{\text{RadFI}}\) be the perturbation operator of
\cref{def:perturbation}, satisfying invariants F1--F5. Let
\(\varepsilon_{\text{SEE}}\) be the canonical SEU operator
parameterized by cross-section \(\sigma\) and flux \(\Phi\). Let
\(X\) count bit-flips per byte of bio payload traversed under
\(\varepsilon_{\text{RadFI}}\) with \(\mathit{prob} = p_R\) and a
uniform I/O workload of duration \(T\), and \(Y\) the corresponding
count under \(\varepsilon_{\text{SEE}}\). Then for every
\(\delta > 0\), there exists \(p_R(\sigma, \Phi)\) such that
\(d_{\text{TV}}(X, Y) \leq \delta\) for \(T\) sufficiently large.

\begin{proof}
Let \(N\) be the number of bytes of bio payload traversed in time
\(T\). Under uniform workload, \(N\) grows linearly:
\(N = r \cdot T\) for I/O rate \(r\).

\emph{Distribution of \(X\).} By F1 and F2, each intercepted bio
that passes the filter and the probability roll produces exactly one
bit-flip at a uniformly sampled byte and bit position. Let \(B\) be
the number of intercepted-and-passed bios in time \(T\). By F4 and
the linearity of the workload,
\(B \sim \text{Binomial}(N / \bar{b}, p_R)\) where
\(\bar{b} = \mathbb{E}[|\text{bio.payload}|]\) is the expected bio
payload size over the I/O workload's size distribution. By the law of large numbers,
\(X = B\) converges in distribution to
\(\text{Poisson}(\lambda_R)\) with
\(\lambda_R = (N / \bar{b}) \cdot p_R\), in the limit
\(N \to \infty\), \(p_R \to 0\), \(N p_R \to \lambda_R\).

\emph{Distribution of \(Y\).} The canonical SEU
model~\cite{dodd2003seu,petersen1996crosssection} gives a Poisson
distribution of bit-flips per byte exposed:
\(Y \sim \text{Poisson}(\lambda_S)\) with
\(\lambda_S = N \cdot \sigma \Phi \cdot 8\) (factor 8 for bits per
byte, under the standard cross-section convention).

\emph{Calibration.} Choose \(p_R = (\bar{b} / N) \cdot \lambda_S\),
i.e.\ \(p_R = \bar{b} \cdot \sigma \Phi \cdot 8\). The dimensional
analysis assumes one device per byte of bio payload (an abuse
standard in software fault injection literature, where the
``cell'' to which \(\sigma\) applies is identified with the byte
on which the flip is observed). This yields \(\lambda_R = \lambda_S\).

\emph{Bounding total variation.} Two Poisson distributions with
equal rate are identically distributed; \(d_{\text{TV}} = 0\). The
convergence of the binomial \(\text{Binomial}(n, p)\) to
\(\text{Poisson}(np)\) is bounded by Le Cam's
theorem~\cite{lecam1960poisson}:
\(d_{\text{TV}}(\text{Binomial}(n, p), \text{Poisson}(np)) \leq n p^2\).
With \(n = N / \bar{b}\) and \(p = p_R\), this gives
\(d_{\text{TV}} \leq (N/\bar{b}) \cdot p_R^2 = \lambda_R \cdot p_R\),
which tends to zero as \(p_R \to 0\) at fixed \(\lambda_R\). For
any \(\delta > 0\), choosing \(p_R \leq \delta / \lambda_R\)
yields \(d_{\text{TV}}(X, Y) \leq \delta\).
\end{proof}

\paragraph{What the proof does not show.}
The proof treats the bit-flip count distribution only. It does not
address: (a) spatial correlation of multi-bit upsets, which is
absent from RadFI by F1; (b) temporal clustering during solar
particle events, which is absent from the uniform-workload
assumption; (c) the joint distribution of flips with the I/O
workload itself. Each is noted as v2 work.
